\worksheettitle
  {Домашняя практика: ответы}
  {\modulemark{M05-H} Ломаная и многоугольник}

\begin{teacherbox}[Критерии]
  Ученик показывает названные элементы на границе и перечисляет вершины
  последовательно. Произвольная перестановка букв не принимается как имя
  многоугольника.
\end{teacherbox}

\section{Ответы}

\begin{enumerate}
  \item Пять вершин; четыре звена; $AB,BC,CD,DE$.
  \item Нужно добавить $NK$; получится четырёхугольник.
  \item Вершины $P,Q,R,S,T$; стороны $PQ,QR,RS,ST,TP$.
    Например, допустимо имя $QRSTP$ или $PTSRQ$.
  \item Для незамкнутой ломаной с шестью вершинами пять звеньев; у замкнутой
    границы с шестью вершинами шесть звеньев. Ошибка --- число точек принято
    за число переходов без проверки замкнутости.
\end{enumerate}

\begin{resultbox}[Маршрут]
  При смешении вершин и звеньев или неверной проверке замкнутости назначьте
  M05-R. Если модель понята, но нарушен порядок букв, дайте короткую
  дополнительную практику. Устойчивый результат открывает M06.
\end{resultbox}
